\subsection{Cost-functorial projection word calculus: certificate equivalence, bidirectional pressure spectrum, and space--time joint phase diagram}\label{subsec:pom-cost-functorial-calculus}
This subsection juxtaposes ``rewritable certificates'' (\(2\)-morphisms of \(\mathbf{PW}\)) with ``cost semantics'' (functors into cost semirings/semifields), so that the following three classes of objects close within a single calculus framework:
(i) strict certification of algorithm equivalence;
(ii) functorialization and invariant/monotone quantities of space--time cost under rewriting;
(iii) the bidirectional pressure spectrum and joint phase diagram \(P(q,u)\) generated jointly by \(\mathrm{PROJ}_u\), the moment-kernel, and \(E_m\) (Schatten spectrum).

\subsubsection{Strict certification of algorithm equivalence: \texorpdfstring{$2$}{2}-morphisms in \texorpdfstring{$\mathbf{PW}$}{PW}}\label{subsec:pom-alg-equivalence-as-2morphism}
\begin{definition}[Certification of algorithm equivalence (exact and \texorpdfstring{$\varepsilon$}{eps}-sound versions)]\label{def:pom-alg-equivalence-2morphism}
In the parametrized projection word \(2\)-category \(\mathbf{PW}\) (Definition \ref{def:pom-projword-2cat}), let \(w,w'\) be \(1\)-morphisms (projection words of composable gates) with the same source and target.
\begin{enumerate}
  \item We say \(w,w'\) are \textbf{exact equivalent} if there exists a \(2\)-morphism certificate
  \[
  \alpha:\ w\Rightarrow w'.
  \]
  \item We say \(w,w'\) are \textbf{exact isomorphic} if there exist \(2\)-morphisms \(\alpha:w\Rightarrow w'\) and \(\beta:w'\Rightarrow w\) such that
  \(
  \beta\circ\alpha=\mathrm{id}_w
  \)
  and
  \(
  \alpha\circ\beta=\mathrm{id}_{w'}
  \);
  denoted \(w\simeq w'\).
  \item Under a given sampling calibration (Definition \ref{def:pom-eps-sound-rewrite}), we say \(w,w'\) are \textbf{\(\varepsilon\)-sound equivalent} if there exists an approximate certificate
  \[
  w\Longrightarrow_{\varepsilon} w'.
  \]
\end{enumerate}
Under the ``auditable calculus'' calibration of this paper, \emph{algorithm equivalence} is no longer stated as ``different implementation details but the same effect''; rather, it is defined as: there exists a composable \(2\)-morphism certificate (or \(\varepsilon\)-sound certificate) in \(\mathbf{PW}\) that rewrites one projection word to another.
\end{definition}

\subsubsection{Conservative semantic functor: collapsing equality tests to natural isomorphism classes}\label{subsec:pom-semantic-conservative-functor}
The rewriting system for projection words does not require global ``completeness'' (i.e., that every semantic equivalence can be derived from the given rules).
From an engineering/audit perspective, what matters more is to construct a set of computable invariants on \emph{bounded complexity slices} that are \emph{conservative} with respect to equivalence classes: identical invariants \(\Rightarrow\) a natural isomorphism must exist.

\begin{definition}[Bounded complexity slice (interface-level)]\label{def:pom-bounded-slice}
Fix parameter upper bounds and a sampling grid
\[
m\le m_{\max},\qquad |G|\le G_{\max},\qquad q\le q_{\max},\qquad \Theta_{\mathrm{grid}}\subset\Theta\ \text{finite}.
\]
Let \(\mathbf{PW}^{\ast}_{\le}\subseteq\mathbf{PW}\) be the full sub-\(2\)-category generated by the minimal derivable fragment
\(\mathbf{PW}^{\ast}=\langle E_m,\mathrm{PROJ}_u,\mathrm{LIFT}_G,\mathrm{TWIST}(\chi),\mathrm{MOM}_q\rangle\)
(Remark \ref{rem:pom-pw-star-fragment}) within the above bounds (objects are the restricted family of interfaces \(\mathcal{I}(m;K;\kappa)\), and \(1\)-morphisms are projection words with correspondingly restricted parameters).
\end{definition}

\begin{remark}[Parameterization by the real line on the unitary slice]\label{rem:pom-unitary-slice-param-by-real-line}
If in subsequent work a parameter axis is restricted to the unitary slice (e.g., taking \(|u|=1\) unit-circle weights under the completion/twisted kernel audit calibration), then the phase variable \(u\in\TT\) can be generated without extra parameters via the Cayley compactification gate of Appendix \ref{app:unit-circle-phase-gate} using the real-line coordinate \(x\in\RR\):
\(
u=\iota(x)=\frac{1+\mathrm{i}x}{1-\mathrm{i}x}
\).
Therefore, ``tuning the phase'' on this slice can be equivalently understood as a choice of coordinate origin, requiring no additional external parameter.
\end{remark}

\begin{definition}[Semantic functor \texorpdfstring{$\mathsf{Sem}$}{Sem} (signature \(\times\) spectrum)]\label{def:pom-sem-functor}
On the slice \(\mathbf{PW}^{\ast}_{\le}\), define an ``auditable semantics'' assignment
\[
\mathsf{Sem}:\mathbf{PW}^{\ast}_{\le}\longrightarrow \mathbf{Sig}_{\le}\times \mathbf{Spec}_{\le}
\]
as follows:
\begin{itemize}
  \item \textbf{Signature part \(\mathbf{Sig}_{\le}\):} Consists of the interface normal form parameters and finite anomaly samples, e.g.,
  \(\bigl(\mathrm{NF}(w),\ G^{\mathrm{BC}}_m(w),\ \mathsf{Anom}_{G}(K;\theta)\vert_{\theta\in\Theta_{\mathrm{grid}}}\bigr)\)
  (where \(G^{\mathrm{BC}}_m(w)\) is the universal abelianization quotient of Definition \ref{def:pom-bc-quotient}: computably extracted from the coupling relation between the folding layer \(E_m\) and the lift axis \(\mathrm{LIFT}_G\) labels in \(w\)).
  (See Proposition \ref{prop:pom-three-gen-interface-normal-form} and Corollary \ref{cor:pom-three-gen-normal-form-decidable}, as well as the structural normal form proposition with \(\mathrm{MOM}\), Proposition \ref{prop:pom-mom-interface-normal-form}).
  \item \textbf{Spectral part \(\mathbf{Spec}_{\le}\):} Consists of computable spectral fingerprints such as the integer moment spectrum and the twisted spectrum, e.g.,
  \(\bigl(r_q(u)\vert_{2\le q\le q_{\max}},\ \rho_\chi(u)/\lambda_1(u)\vert_{\chi\in\widehat G,\,\theta\in\Theta_{\mathrm{grid}}}\bigr)\)
  and their derived pressure/dual budget objects (see \S\ref{subsec:pom-quantile-budget-law} and \S\ref{subsec:pom-spectral-resource-type}).
\end{itemize}
In the implementation of this paper, each coordinate of \(\mathsf{Sem}\) is extracted by the closed-loop interface ``matrix family \(\Rightarrow\zeta\Rightarrow\) finite part/spectral radius/moment spectrum'' (\S\ref{subsec:pom-implementable-interfaces}).
\end{definition}

\begin{theorem}[Conservativity (bounded slice): \texorpdfstring{$\mathsf{Sem}$}{Sem} reflects natural isomorphisms]\label{thm:pom-sem-conservative}
On the slice of Definition \ref{def:pom-bounded-slice}, \(\mathsf{Sem}\) is conservative: for any projection words \(w_1,w_2\in\mathbf{PW}^{\ast}_{\le}\) with the same source and target, if
\[
\mathsf{Sem}(w_1)=\mathsf{Sem}(w_2),
\]
then there must exist a natural isomorphism \(w_1\simeq w_2\) (Definition \ref{def:pom-alg-equivalence-2morphism}).
\end{theorem}

\begin{remark}[Implication: no global completeness needed, only slice conservativity]\label{rem:pom-sem-conservative-meaning}
Theorem \ref{thm:pom-sem-conservative} promotes ``algorithm equivalence'' from ``same final kernel fingerprint'' to a decidable version of ``same natural isomorphism class'':
within the bounded slice, there is no need to prove that the rewriting system is complete for all semantic equivalences; it suffices to ensure that the semantic output \(\mathsf{Sem}\) is fine enough to reflect isomorphisms on that slice.
This aligns with the completion/audit strategy of this paper: under fixed \(m,G,q\) and a finite parameter grid, verifying ``normal form consistency + residual class consistency'' through critical pair enumeration and regression scripts (Remark \ref{rem:pom-threegen-kb-completion-audit}, Remark \ref{rem:pom-fourgen-kb-completion-audit}) suffices to collapse the equivalence test to a reproducible finite output comparison.
\end{remark}

\begin{remark}[Finite signaturization: unifying ``equivalence + optimization'' as the same set of auditable outputs]\label{rem:pom-finite-signature-sigmam}
Within the bounded slice, the key coordinates of \(\mathsf{Sem}\) can be compressed into a finite signature family closer to a ``calculus-level acceptance criterion'':
for any projection word \(w\) (whose resolution, lift group, and temperature/potential parameters are all bounded by the slice constraints), define
\[
\Sigma_m(w)
:=
\Bigl(
G^{\mathrm{BC}}_m(w),\
\{\mathsf{Anom}_{G}(K;\theta)\}_{\theta\in\Theta_{\mathrm{grid}}},\
\{r_q(u)\}_{(q,u)\in\mathcal{G}_{\mathrm{grid}}}
\Bigr),
\]
where \(G^{\mathrm{BC}}_m(w)\) is the universal abelianization quotient (Definition \ref{def:pom-bc-quotient}), and \(\Theta_{\mathrm{grid}}\) and \(\mathcal{G}_{\mathrm{grid}}\subseteq\{2,\dots,q_{\max}\}\times(0,1]\) are finite test grids.
Then ``equality testing'' and ``optimization'' can be uniformly understood as: among candidate realizations satisfying \(\Sigma_m\) consistency, take the Pareto minimum in \(\mathbb{T}^{(2)}\) via the cost functor (see Proposition \ref{prop:pom-cost-optimization-decidable}).
\end{remark}

\begin{remark}[Computable closed loop: rewriting system \texorpdfstring{$+$}{+} semantic verification system]\label{rem:pom-rewrite-semantic-check-loop}
The organization of this section can be compressed into an executable ``calculus closed loop'':
\begin{enumerate}
  \item \textbf{Syntactic layer (certified rewriting):}
  Perform normalization and rewriting in \(\mathbf{PW}\), outputting the interface normal form \(\mathrm{NF}(w)\) and its \(2\)-morphism certificate (Definition \ref{def:pom-alg-equivalence-2morphism}).
  When approximate compilation is allowed, the \(\varepsilon\)-sound certificate carries a risk budget (Definition \ref{def:pom-eps-sound-rewrite}), and the minimal budget distance \(d_{\mathcal I}\) internalizes the budget as an intrinsic quantity on homomorphism classes (Definition \ref{def:pom-eps-distance}; Proposition \ref{prop:pom-eps-distance-lawvere}).
  \item \textbf{Semantic layer (finite signature/spectral fingerprint):}
  The conservative semantic functor \(\mathsf{Sem}\) sends words into the computable output domain of ``signature \(\times\) spectrum'' (Definition \ref{def:pom-sem-functor}), and verification on bounded slices uses the finite signature \(\Sigma_m(w)\) (Remark \ref{rem:pom-finite-signature-sigmam}).
  \item \textbf{Verification layer (equality/consistency):}
  Under the slice calibration, the consistency of \(\mathsf{Sem}\) reflects natural isomorphisms (Theorem \ref{thm:pom-sem-conservative}), so ``lexical normal form consistency + residual class consistency + finite spectral fingerprint consistency'' yields a terminating equality-testing procedure;
  for approximate compilation, \(d_{\mathcal I}(w_1,w_2)\le\varepsilon\) serves as the intrinsic criterion at the ``allowable error'' calibration.
  Moreover, the ``infinite registers'' of the congruence/character axis do not obstruct termination: any finite projection word touches only finitely many modular supports, and verification can be reduced to exhaustive enumeration on some finite layer \(\widehat{G_{m_\ast}}\) (Proposition \ref{prop:pom-finite-modulus-stability}).
\end{enumerate}
\end{remark}

\subsubsection{Cost semiring and cost functor: from \texorpdfstring{$\mathbf{PW}$}{PW} to the resource preorder}\label{subsec:pom-cost-functor}
\paragraph{Cost semiring}
Let \(\overline{\RR}_{\ge 0}:=\RR_{\ge 0}\cup\{+\infty\}\), and set
$$
\mathbb{T}:=\bigl(\overline{\RR}_{\ge 0},\ \oplus,\ \otimes,\ \mathsf 0,\ \mathsf 1\bigr),
\qquad
a\oplus b:=\min(a,b),\quad a\otimes b:=a+b,
\qquad
\mathsf 0:=+\infty,\ \mathsf 1:=0.
$$
Then \(\mathbb{T}\) is the idempotent commutative semiring (tropical semiring). For the joint space--time resource, take the product semiring
$$
\mathbb{T}^{(2)}:=\mathbb{T}\times \mathbb{T},
\qquad
(s,t)\oplus(s',t'):=(\min(s,s'),\min(t,t')),
\quad
(s,t)\otimes(s',t'):=(s+s',t+t').
$$
\(\mathbb{T}^{(2)}\) carries the pointwise preorder \((s,t)\preceq(s',t')\iff s\le s'\ \wedge\ t\le t'\).

\begin{definition}[Resource preorder \(2\)-category and cost functor]\label{def:pom-cost-functor}
Define a thin \(2\)-category \(\mathbf{Res}\) (resources) as follows:
\begin{itemize}
  \item \textbf{Objects:} The same interface objects as in \(\mathbf{PW}\) (Definition \ref{def:pom-projword-2cat}, \(\mathcal{I}(m;K;\kappa)\)).
  \item \textbf{\(1\)-morphisms (costs):} For any pair of objects \((A,B)\), let \(\mathrm{Hom}_1(A,B):=\mathbb{T}^{(2)}\).
  \item \textbf{\(2\)-morphisms (comparisons):} For any \(c,c'\in\mathbb{T}^{(2)}\), there exists a unique \(2\)-morphism
  \(
  c\Rightarrow c'
  \)
  if and only if \(c'\preceq c\).
  \item \textbf{Composition and identity:} Given by \(\otimes\):
  \[
  c_2\circ c_1:=c_2\otimes c_1,
  \qquad
  \mathrm{id}_A:=(0,0).
  \]
\end{itemize}
A (weak) \(2\)-functor \(\mathcal{C}:\mathbf{PW}\to \mathbf{Res}\) is called a \textbf{cost functor} if on \(1\)-morphisms it satisfies
\[
\mathcal{C}(w_2\circ w_1)=\mathcal{C}(w_2)\otimes \mathcal{C}(w_1),
\qquad
\mathcal{C}(\mathrm{id}_A)=(0,0),
\]
and for any rewriting certificate \(\alpha:w\Rightarrow w'\),
\[
\mathcal{C}(w')\preceq \mathcal{C}(w).
\]
If one wishes to regard a certain class of certificates as strict \emph{cost invariants} (invariant calibration), the additional requirement is that \(\mathcal{C}(w)=\mathcal{C}(w')\) holds identically for certificates of that class.
\end{definition}

\begin{remark}[Stratification of invariants and monotone quantities]\label{rem:pom-cost-invariant-vs-monotone}
Definition \ref{def:pom-cost-functor} allows ``rewriting'' to be divided into two types: one type is strict equivalence (which should be invariant in both semantics and cost), and the other is compilation-type rewriting (which is allowed to change cost but is required to be monotone under the chosen preorder).
The typical usage in this paper is: exact equality testing is delegated to the \(\Sigma/\zeta/\)finite-part fingerprint functors (Definition \ref{def:pom-sigma-functor} and Remark \ref{rem:pom-zeta-normal-form}), while ``space-for-time/time-for-space'' compilation laws are written as monotone certificates in \(\mathbb{T}^{(2)}\) (e.g., the failure probability ledger and prime-register quota of \((\mathrm{RCRT}_\varepsilon)\)).
\end{remark}

\begin{definition}[Cost-enriched projection word calculus \texorpdfstring{$\mathbf{PW}^{\mathrm{cost}}$}{PWcost}]\label{def:pom-pw-cost}
Given a cost functor \(\mathcal{C}:\mathbf{PW}\to\mathbf{Res}\) (Definition \ref{def:pom-cost-functor}), the pair
\[
\mathbf{PW}^{\mathrm{cost}}:=\bigl(\mathbf{PW},\mathcal{C}\bigr)
\]
is called the \textbf{cost-enriched projection word calculus}.
In this structure, ``semantic equivalence'' is given by \(2\)-isomorphisms (or \(\varepsilon\)-sound certificates) in \(\mathbf{PW}\), while ``compilation optimization'' is understood as: within a fixed semantic equivalence class, finding the Pareto-minimal cost realization under the resource preorder \((\mathbb{T}^{(2)},\preceq)\).
\end{definition}

\begin{proposition}[Compilation optimization is decidable on bounded slices (finite frontier)]\label{prop:pom-cost-optimization-decidable}
On the bounded slice of Definition \ref{def:pom-bounded-slice}, fix an auditable fragment \(\mathbf{PW}^{\ast}_{\le}\subseteq\mathbf{PW}\) and a computable cost functor \(\mathcal{C}\).
Then for any projection word \(w\in\mathbf{PW}^{\ast}_{\le}\), its semantic equivalence class
\[
[w]:=\{\,w'\in\mathbf{PW}^{\ast}_{\le}:\ w'\simeq w\,\}
\]
has the following properties for its image in the cost domain
\(
\mathcal{C}([w])\subseteq \mathbb{T}^{(2)}
\):
\begin{enumerate}
  \item \textbf{(Finite frontier)} The Pareto-minimal element set of \(\mathcal{C}([w])\) is a finite set;
  \item \textbf{(Decidable/constructible)} The Pareto frontier can be computed by finite enumeration: within the parameter grid and complexity upper bounds given by the slice, enumerate candidate interface normal forms and use \(\mathsf{Sem}\) (Theorem \ref{thm:pom-sem-conservative}) for equality-testing binning, then take the \(\preceq\)-minimal elements within each bin.
\end{enumerate}
\end{proposition}
\begin{proof}[Proof sketch]
On the bounded slice, generator parameters are drawn from a finite grid, and the key fragments of this paper all provide terminating normalization strategies that reduce any word to a finite-length family of interface normal forms (e.g., three/four-generator normal forms, Propositions \ref{prop:pom-three-gen-interface-normal-form} and \ref{prop:pom-mom-interface-normal-form}).
Therefore, the interface normal form parameter combinations that can arise within this slice are only finitely many; consequently the semantic signature bins also number only finitely many.
For each bin, the cost functor \(\mathcal{C}\) is computable under generators and composition, so the \(\preceq\)-minimal elements can be directly taken from the finite candidate set to obtain the Pareto frontier.
\end{proof}

\subsubsection{Tropicalization of the temperature gate: from PF dominant spectrum to min-plus eigenvalue}\label{subsec:pom-proj-tropicalization}
\begin{definition}[Tropical cost of the temperature gate]\label{def:pom-proj-tropical-cost}
Following the matrix realization of the temperature-parametrized order projection \(\mathrm{PROJ}_u\) (Definition \ref{def:pom-proj-temp}), for the weighted adjacency matrix with potential \(M_E(u)\) define its \textbf{tropical cost} as
$$
c_{\mathrm{trop}}(E)
:=\lim_{u\downarrow 0}\frac{\log \rho\!\bigl(M_E(u)\bigr)}{\log u},
$$
whenever the limit exists.
\end{definition}

\begin{proposition}[Tropical limit = minimum mean cycle energy]\label{prop:pom-proj-tropical-minmean}
Under the setting of a finite graph kernel (finite-state SFT) with nonnegative integer edge potentials \(E\), the limit in Definition \ref{def:pom-proj-tropical-cost} exists and satisfies
$$
\boxed{\;
c_{\mathrm{trop}}(E)
=
\min_{\gamma\ \mathrm{cycle}}\ \frac{E(\gamma)}{|\gamma|}\;},
$$
where \(\gamma\) ranges over directed cycles of the graph, \(E(\gamma)\) is the potential sum along the cycle, and \(|\gamma|\) is the cycle length.
\end{proposition}
\begin{proof}[Proof sketch (combinatorial certificate)]
Set \(u=e^{-\beta}\) (\(\beta\to+\infty\) is equivalent to \(u\downarrow 0\)).
For edge potentials on a finite graph, the zero-temperature limit concentrates the equilibrium onto the ground-state set of minimum average cost (Remark \ref{rem:pom-order-as-zero-temp}), and the minimum average cost equals the min-mean-cycle value on the graph, computable by Karp's algorithm (same remark). On the other hand, \(\rho(M_E(e^{-\beta}))\) has a slope limit in its logarithm as \(\beta\to\infty\), and this slope is precisely the above minimum average cost; converting back to the variable \(u\) yields
\(
\log\rho(M_E(u))\sim c_{\mathrm{trop}}(E)\log u
\)
and the conclusion follows.
\end{proof}

\begin{remark}[The same gate family unifies average and worst case]\label{rem:pom-proj-avg-worst-unified}
\(\mathrm{PROJ}_u\) at \(u\in(0,1)\) gives a thermodynamic (average) readout; Proposition \ref{prop:pom-proj-tropical-minmean} shows that at \(u\to 0\) it automatically degenerates to the min-plus readout of the ``most energy-efficient certificate'' (worst case/hard constraint). Therefore, once the lexical-level certificate rewriting is completed in \(\mathbf{PW}\) (Definition \ref{def:pom-alg-equivalence-2morphism}), the same certificate can be pushed by \(u\) to comparison semantics at both the ``average'' and ``worst case'' ends.
\end{remark}

\subsubsection{Bidirectional pressure spectrum of positive and negative moments: moment-kernel \texorpdfstring{$+$}{+} Schatten spectrum}\label{subsec:pom-bidirectional-pressure}
\begin{definition}[Bidirectional moments and bidirectional pressure]\label{def:pom-bidirectional-moments}
For each resolution \(m\), denote the fiber multiplicity \(d_m(x)=|\Fold_m^{-1}(x)|\) (\(x\in X_m\)).
For any \(t>0\), define the \textbf{positive moment} and \textbf{negative moment}:
$$
S^{+}_{t}(m):=\sum_{x\in X_m} d_m(x)^{t},
\qquad
S^{-}_{t}(m):=\sum_{x\in X_m} d_m(x)^{-t}.
$$
And define the corresponding exponential-scale (upper) pressure as
$$
\tau^{+}(t):=\limsup_{m\to\infty}\frac{1}{m}\log S^{+}_{t}(m),
\qquad
\widetilde P(t):=\tau^{-}(t):=\limsup_{m\to\infty}\frac{1}{m}\log S^{-}_{t}(m).
$$
\end{definition}

\begin{remark}[Spectral operator realization of space loss: \texorpdfstring{$\mathcal{F}_m=(E_mE_m^\ast)^{-1}$}{Fm=(EmEm*)^{-1}}]\label{rem:pom-fiber-loss-as-spectral-operator}
By Appendix Definition \ref{def:fold-fiber-multiplicity-operator} and Proposition \ref{prop:fold-fiber-multiplicity-trace-moments}, the fiber multiplicity can be promoted to a positive diagonal spectral operator \(\mathcal{F}_m\) on \(\ell^2(X_m)\) (whose spectrum is \(\{d_m(x)\}_{x\in X_m}\)).
Under this calibration, the bidirectional moments are themselves the trace moments of \(\mathcal{F}_m\):
$$
\boxed{\;
S^{+}_{t}(m)=\Tr(\mathcal{F}_m^{t}),\qquad
S^{-}_{t}(m)=\Tr(\mathcal{F}_m^{-t})\; }.
$$
Therefore the pressure functions \(\tau^{\pm}(t)\) can be understood as the thermodynamic limits of the trace moments of this spectral operator; the PF dominant spectrum of the moment-kernel and the Schatten spectrum of \(E_m\) correspond respectively to the positive-side and negative-side trace moment interfaces of \(\mathcal{F}_m\).
\end{remark}

\begin{remark}[Negative moments are directly read from the Schatten spectrum of \texorpdfstring{$E_m$}{Em}]\label{rem:pom-negative-moment-from-schatten}
By the Appendix operator algebra calibration (Corollary \ref{cor:fold-conditional-expectation-schatten}), when \(t\ge \frac12\) (i.e., \(2t\ge 1\)), the exact identity holds:
$$
\boxed{\;
S^{-}_{t}(m)=\sum_{x\in X_m}d_m(x)^{-t}=\|E_m\|_{S_{2t}}^{2t}\;}.
$$
Therefore the negative moment spectrum is a directly computable semantic output of the \(\mathbf{PW}\) generator \(E_m\): it requires no additional moment-kernel compilation.
\end{remark}

\begin{proposition}[Positive--negative moment coupling inequality (space spectral uncertainty lower bound)]\label{prop:pom-positive-negative-coupling}
For any \(t>0\) and any \(m\),
$$
\boxed{\;
S^{+}_{t}(m)\,S^{-}_{t}(m)\ \ge\ |X_m|^2\; }.
$$
Hence on the exponential scale,
$$
\tau^{+}(t)+\widetilde P(t)\ \ge\ 2\liminf_{m\to\infty}\frac{1}{m}\log|X_m|
=2\log\varphi.
$$
\end{proposition}
\begin{proof}
Set \(a_x=d_m(x)^{t/2}\), \(b_x=d_m(x)^{-t/2}\), so that \(\sum_x a_xb_x=|X_m|\).
By Cauchy--Schwarz,
\(
|X_m|^2\le (\sum_x a_x^2)(\sum_x b_x^2)=S_t^+(m)S_t^-(m)
\).
Taking logarithms, dividing by \(m\), and noting \(|X_m|=F_{m+2}=\Theta(\varphi^m)\) yields the second inequality.
\end{proof}

\begin{remark}[Cross-audit inequality for the two pipelines]\label{rem:pom-positive-negative-cross-audit}
For any \(t\ge \frac12\), by Remark \ref{rem:pom-negative-moment-from-schatten} one can rewrite Proposition \ref{prop:pom-positive-negative-coupling} as
$$
\boxed{\;
S^{+}_{t}(m)\,\|E_m\|_{S_{2t}}^{2t}\ \ge\ |X_m|^2\; }.
$$
Here \(S_t^{+}(m)\) is computed via the moment-kernel/residue DP pipeline (positive moment pipeline), while \(\|E_m\|_{S_{2t}}\) can be directly read from the fiber histogram \(d_m(\cdot)\) or the spectral fingerprint of Table \ref{tab:fold_conditional_expectation_singular_spectrum} (operator pipeline).
Therefore this inequality provides a deterministic consistency constraint for two mutually independent audit subsystems.
\end{remark}

\begin{corollary}[Fiber dispersion index: arithmetic/harmonic mean ratio]\label{cor:pom-fiber-dispersion-index}
Let \(\overline d_m:=2^m/|X_m|\) be the arithmetic mean of the fiber size under uniform \(X_m\) sampling, and let
\[
d_m^{\mathrm{harm}}
:=
\frac{|X_m|}{\sum_{x\in X_m} d_m(x)^{-1}}
=
\frac{|X_m|}{\|E_m\|_{HS}^2}
\]
be the corresponding harmonic mean (Corollary \ref{cor:fold-conditional-expectation-schatten}).
Define the \textbf{fiber dispersion index}
\[
\boxed{\;
\mathfrak{D}_m
:=
\frac{\overline d_m}{d_m^{\mathrm{harm}}}
=
\frac{2^m}{|X_m|^2}\sum_{x\in X_m}\frac{1}{d_m(x)}
=
\frac{2^m}{|X_m|^2}\|E_m\|_{HS}^2\; }.
\]
Then for all \(m\ge 1\) the deterministic lower bound holds:
\[
\boxed{\;\mathfrak{D}_m\ge 1\;}
\]
with equality if and only if all fiber sizes \(d_m(x)\) are everywhere equal (completely uniform fibers).
\end{corollary}
\begin{proof}
By Proposition \ref{prop:pom-positive-negative-coupling} with \(t=1\), noting that \(S_1^+(m)=\sum_x d_m(x)=|\Omega_m|=2^m\), we obtain
\[
2^m\sum_{x\in X_m}\frac{1}{d_m(x)}\ \ge\ |X_m|^2,
\]
and dividing both sides by \(|X_m|^2\) gives \(\mathfrak{D}_m\ge 1\).
The equality condition follows from the necessary and sufficient condition for equality in Cauchy--Schwarz: \(\bigl(d_m(x)^{1/2}\bigr)_x\) and \(\bigl(d_m(x)^{-1/2}\bigr)_x\) are proportional if and only if \(d_m(x)\) is constant.
\end{proof}
\input{sections/generated/tab_fold_fiber_dispersion_index}

\begin{remark}[Heterogeneity trend within the audit window]\label{rem:pom-fiber-dispersion-trend}
Table \ref{tab:fold_fiber_dispersion_index} provides reproducible output of \(\mathfrak{D}_m\) for several values of \(m\); within the audit window \(m=8,10,\dots,24\), \(\mathfrak{D}_m\) slowly increases from approximately \(1.2367\) to approximately \(1.3781\), indicating that the heterogeneity of folding fibers grows with increasing resolution.
This scalar differs from both the maximum fiber statistic and the positive moment spectrum \(S_t^+(m)\);
it measures the heterogeneity between ``typical fibers'' and ``sharp fibers'' in the form of a mean ratio, subject to the hard constraint of the lower bound \(\mathfrak{D}_m\ge 1\).
\end{remark}

\begin{remark}[Bridge to the positive moment pressure \texorpdfstring{$P(q)$}{P(q)}]\label{rem:pom-negative-pressure-bridge-to-positive}
When \(t=q\in\ZZ_{\ge 2}\), \(S_q^+(m)=S_q(m)\) is given by the moment-kernel as \(\tau^{+}(q)=\log r_q\) (Theorem \ref{thm:pom-collision-kernel-ak}).
Therefore Proposition \ref{prop:pom-positive-negative-coupling} provides a lower bound that hard-couples the two semantic chains:
$$
\boxed{\;
\widetilde P(q)\ \ge\ 2\log\varphi-\log r_q\; }.
$$
This translates the calibration that ``space loss is not a number but a spectrum'' into a verifiable inequality: the heavier the positive moment tail (large fibers), the more the negative moment tail (small fibers/sharp fibers) is forced to increase.
\end{remark}

\paragraph{Interface-level output of the negative moment--Schatten left-tail channel (audit interface)}\label{par:pom-left-tail-schatten-terminal}
Beyond the positive moment collision spectrum \(\{S_q(m)\}_{q\ge 2}\) (moment-kernel pipeline) of the main text, Appendix \ref{app:op-algebra} uses the singular value spectrum of the folding conditional expectation \(E_m\) as its direct interface to provide a \textbf{negative moment--Schatten left-tail channel} orthogonal to the former (Remark \ref{rem:pom-negative-moment-from-schatten}).
This channel translates the empirical description of ``low-multiplicity layers'' into a reproducible certificate system: all quantities can be directly verified from Table \ref{tab:fold_conditional_expectation_singular_spectrum} and its derived Tables \ref{tab:fold_conditional_expectation_left_tail_bases} and \ref{tab:fold_conditional_expectation_effective_support}.

\medskip\noindent\textbf{The left-tail free energy density constant family exhibits a stable base within the audit window.}
Let
\[
T_s(m):=S_s^{-}(m)=\sum_{x\in X_m} d_m(x)^{-s},
\qquad
t_s:=\limsup_{m\to\infty}\frac{1}{m}\log T_s(m)=\tau^{-}(s).
\]
On the even grid \(m=8,10,\dots,24\) covered by the audit window, take the empirical base
\[
b_s(m):=\Bigl(\frac{T_s(m)}{T_s(m-2)}\Bigr)^{1/2}\qquad(m\ge 10),
\]
then Table \ref{tab:fold_conditional_expectation_left_tail_bases} shows rapid stabilization at \(s=\tfrac12,1,2\); at the audit point \(m=24\),
\[
b_{1/2}(24)\approx 1.458200,\qquad b_1(24)\approx 1.316686,\qquad b_2(24)\approx 1.092366,
\]
giving the candidate exponential rates (natural logarithm)
\[
t_{1/2}\approx 0.3772,\qquad t_1\approx 0.2751,\qquad t_2\approx 0.0884.
\]

\medskip\noindent\textbf{The negative moment channel requires no moment-kernel compilation; its information is directly read from the Schatten spectrum of \(E_m\).}
Under the operator interface \(\mathcal{F}_m=(E_mE_m^\ast)^{-1}\),
\[
T_s(m)=\Tr(\mathcal{F}_m^{-s}),
\qquad
\text{and when }s\ge \tfrac12\text{, }T_s(m)=\|E_m\|_{S_{2s}}^{2s}
\]
(Appendix Proposition \ref{prop:fold-fiber-multiplicity-trace-moments}).
Therefore the left-tail negative moment spectrum is a direct semantic output of the generator \(E_m\), independent of any additional construction of higher-order collision kernels.

\medskip\noindent\textbf{Two types of effective support scale give a two-scale condensation of singular value energy, with stable separation emerging at \(m=24\).}
Let the singular value energy distribution be \(p_x\propto d_m(x)^{-1}\), and define
\[
N_2(m):=\frac{1}{\sum_x p_x^2},
\qquad
N_{1/2}(m):=\Bigl(\sum_x\sqrt{p_x}\Bigr)^2
\]
(Table \ref{tab:fold_conditional_expectation_effective_support} is read in closed form from Schatten norms).
Within the audit window \(m=8,10,\dots,24\), \(\frac{N_2(m)}{|X_m|}\) monotonically decreases, while \(\frac{N_{1/2}(m)}{|X_m|}\) remains in an interval close to \(1\) throughout (Table \ref{tab:fold_conditional_expectation_effective_support}).
At the audit point \(m=24\),
\[
\frac{N_2(24)}{|X_{24}|}\approx 0.455665,\qquad
\frac{N_{1/2}(24)}{|X_{24}|}\approx 0.895533,\qquad
\frac{N_{1/2}(24)}{N_2(24)}\approx 1.9653.
\]
This indicates: in the R\'enyi-2 (collision) sense, the energy covers only about \(0.46|X_m|\) effective scale, while in the R\'enyi-\(\tfrac12\) (Bhattacharyya) sense it still covers about \(0.90|X_m|\) effective scale; this two-scale separation constitutes a reproducible spectral fingerprint.

\medskip\noindent\textbf{The low-multiplicity layer admits a parameter-free left-tail certificate, with numerically quantified lower bounds available at \(m=24\).}
Let \(X\sim \mathrm{Unif}(X_m)\).
For any \(s>0\), set
\[
d_{-}(m,s):=\Bigl(\frac{2|X_m|}{T_s(m)}\Bigr)^{1/s}.
\]
Then by the negative moment Paley--Zygmund certificate (Appendix Proposition \ref{prop:fold-negative-moment-left-tail-certificate}),
\[
\mathbb{P}\bigl(d_m(X)\le d_{-}(m,s)\bigr)\ \ge\ \frac14\,\frac{T_s(m)^2}{|X_m|\,T_{2s}(m)}.
\]
At the reproducible Schatten output at \(m=24\) (Tables \ref{tab:fold_conditional_expectation_singular_spectrum} and \ref{tab:fold_conditional_expectation_effective_support}), one directly obtains two numerical-level certificates:
\[
\mathbb{P}\bigl(d_{24}(X)\le 200.57\bigr)\ge 11.39\% \quad(s=1),
\qquad
\mathbb{P}\bigl(d_{24}(X)\le 447.94\bigr)\ge 22.39\% \quad(s=\tfrac12).
\]
The required inputs are only \(\|E_m\|_{S_1}\), \(\|E_m\|_{HS}^2\), and \(\|E_m\|_{S_4}\) from Table \ref{tab:fold_conditional_expectation_singular_spectrum} (yielding \(T_{1/2},T_1,T_2\)), together with the effective support scale readout from Table \ref{tab:fold_conditional_expectation_effective_support}.

\medskip\noindent\textbf{The space spectral uncertainty lower bound binds right-tail compression and left-tail thinning as a mutually exclusive constraint.}
Proposition \ref{prop:pom-positive-negative-coupling} gives
\(
S_t^+(m)\,S_t^-(m)\ge |X_m|^2
\)
together with the exponential-scale coupling
\(
\tau^{+}(t)+\tau^{-}(t)\ge 2\log\varphi
\).
Therefore, any strategy that increases the positive moment pressure to strengthen right-tail certificates must necessarily push up the left-tail pressure by the same order of magnitude, thereby amplifying the exponential scale of the low-multiplicity layer; this constraint is unavoidable at the certificate level.

\medskip\noindent\textbf{Two-tail certificates close into a quantile pinching band on the same audit interface and can be juxtaposed with the prime-register budget.}
Juxtaposing the right-tail positive moment Markov/Chernoff upper bound with the left-tail negative moment Paley--Zygmund lower bound yields explicit threshold functions:
for any integer \(q\ge 1\), \(\varepsilon\in(0,1)\), set
\[
d_{+}(m,q,\varepsilon):=\Bigl(\frac{S_q(m)}{\varepsilon\,|X_m|}\Bigr)^{1/q}.
\]
Then under \(X\sim \mathrm{Unif}(X_m)\),
\[
\mathbb{P}\bigl(d_m(X)\le d_{-}(m,s)\bigr)\ \text{has an explicit lower bound (from the negative moment Paley--Zygmund certificate)},
\qquad
\mathbb{P}\bigl(d_m(X)\ge d_{+}(m,q,\varepsilon)\bigr)\le \varepsilon.
\]
This means that at the same resolution \(m\), the left-tail proportion lower bound (Schatten negative moment interface) and the right-tail upper bound (moment-kernel positive moment interface) can be independently audited and juxtaposed, thereby providing an additional control surface for the prime-register tail budget constraint that is orthogonal to the low-multiplicity structure yet closed at the data interface level.

\subsubsection{RCRT spatialization pressure reduction law: fiber refinement by residue necessarily disperses collisions}\label{subsec:pom-rcrt-pressure-reduction}
\begin{definition}[Collision moment after residue refinement]\label{def:pom-residue-refined-moment}
Fix the modulus product \(P\) (prime-register capacity), and refine each fiber by residue class: for each \(x\in X_m\), let
$$
d_m(x)=\sum_{i=1}^{P} d_{m,i}(x),
$$
where \(d_{m,i}(x)\) denotes the number of microstates falling into the \(i\)-th residue bucket (allowed to be \(0\)).
For integer \(q>1\), define the refined \(q\)-moment
$$
S_q^{(P)}(m):=\sum_{x\in X_m}\sum_{i=1}^{P} d_{m,i}(x)^q.
$$
\end{definition}

\begin{lemma}[Upper and lower bounds for the refined moment (Jensen/convexity)]\label{lem:pom-residue-refinement-jensen}
For any \(q>1\) and any nonnegative sequence \((d_i)_{i=1}^P\) with \(\sum_i d_i=d\),
$$
\boxed{\;
P^{1-q}d^q\ \le\ \sum_{i=1}^{P}d_i^q\ \le\ d^q\; }.
$$
\end{lemma}
\begin{proof}
The right-hand side follows from \(\sum_i d_i^q\le (\sum_i d_i)^q=d^q\) (superadditivity of \(t\mapsto t^q\) for \(q>1\)).
The left-hand side follows from convexity (Jensen): \(\frac{1}{P}\sum_i d_i^q\ge (\frac{1}{P}\sum_i d_i)^q=(d/P)^q\).
\end{proof}

\begin{corollary}[Explicit pressure reduction law]\label{cor:pom-rcrt-pressure-reduction-law}
Under the setting of Definition \ref{def:pom-residue-refined-moment}, for any \(q>1\) and any \(m\),
$$
P^{1-q}S_q(m)\ \le\ S_q^{(P)}(m)\ \le\ S_q(m).
$$
Therefore, for the fixed-\(m\) finite-scale pressure
$$
P_m^{(P)}(q):=\frac{1}{m}\log S_q^{(P)}(m)-q\log 2
$$
we have
$$
P_m(q)-\frac{q-1}{m}\log P\ \le\ P_m^{(P)}(q)\ \le\ P_m(q),
$$
where \(P_m(q):=\frac{1}{m}\log S_q(m)-q\log 2\).
If we choose \(P=P_m\) varying with \(m\) such that
\(
\sigma:=\lim_{m\to\infty}\frac{1}{m}\log P_m
\)
exists (e.g., exponentially growing prime-register capacity), then setting
\[
P^{(P_\bullet)}(q):=\limsup_{m\to\infty} P_m^{(P_m)}(q),
\]
and under the calibration \(P_m(q)\to P(q)\), we obtain the computable reduction band
$$
\boxed{\;
P(q)- (q-1)\sigma\ \le\ P^{(P_\bullet)}(q)\ \le\ P(q)\; }.
$$
\end{corollary}
\begin{proof}
Apply Lemma \ref{lem:pom-residue-refinement-jensen} for each \(x\) and sum over \(x\) to get the first inequality; the rest follows by taking logarithms and dividing by \(m\).
\end{proof}

\begin{remark}[Implication: spatialization is not merely a gate replacement but changes the collision pressure itself]\label{rem:pom-rcrt-pressure-meaning}
When (RCRT) (Definition \ref{def:pom-order-elimination-rewrite}) is used for order-gate elimination, the residue axis of the prime-register not only provides a ``uniquely reconstructible'' semantic channel; the above inequality shows that the same refinement necessarily disperses collisions at the statistical level, thereby suppressing the positive moment pressure and the multifractal upper envelope (Proposition \ref{prop:pom-fold-multifractal-envelope}).
This gives ``space-for-time'' a verifiable spectral-side transformation law: purchasing residue dimensions with \(\log P\) produces at least \((q-1)\log P\) worth of exponential reduction space on the \(q\)-moment.
\end{remark}

\begin{corollary}[Self-consistent square-root quota law: the same prime-register simultaneously serves as carrier and disperser]\label{cor:pom-rcrt-self-consistent-sqrt-budget}
Under the finite-scale \(q\)-certificate calibration of Corollary \ref{cor:pom-prime-register-q-quota-family}, for integer \(q\ge 2\) define
\[
B_{m,q}^{\mathrm{fin}}(\varepsilon)
:=\varepsilon^{-1/(q-1)}\left(\frac{S_q(m)}{2^m}\right)^{\!1/(q-1)}.
\]
Under the residue refinement of Definition \ref{def:pom-residue-refined-moment}, define the refined certificate correspondingly as
\[
B_{m,q}^{(P),\mathrm{fin}}(\varepsilon)
:=\varepsilon^{-1/(q-1)}\left(\frac{S_q^{(P)}(m)}{2^m}\right)^{\!1/(q-1)}.
\]
Then by Corollary \ref{cor:pom-rcrt-pressure-reduction-law}, the pinching immediately follows:
\[
\boxed{\;
\frac{1}{P}\,B_{m,q}^{\mathrm{fin}}(\varepsilon)
\ \le\
B_{m,q}^{(P),\mathrm{fin}}(\varepsilon)
\ \le\
B_{m,q}^{\mathrm{fin}}(\varepsilon)\;}
\]
(i.e., the same residue refinement gives at most a \(P\)-fold and at least no improvement on the \(q\)-moment tail certificate).
Therefore, for the same prime-register capacity variable \(P\):
\begin{itemize}
  \item \textbf{(Conservative threshold; not counting on dispersal improvement)} If \(P>2B_{m,q}^{\mathrm{fin}}(\varepsilon)\), then by Corollary \ref{cor:pom-prime-register-q-quota-family} and Proposition \ref{prop:pom-rcrt-epsilon}, a sufficient condition for the \(\varepsilon\)-sound compilation of \((\mathrm{RCRT}_\varepsilon)\) is obtained.
  \item \textbf{(Square-root frontier when Jensen bound is nearly tight)} If the residue buckets are statistically near-uniform, so that \(S_q^{(P)}(m)\approx P^{1-q}S_q(m)\) (Jensen lower bound nearly attained), then \(B_{m,q}^{(P),\mathrm{fin}}(\varepsilon)\approx B_{m,q}^{\mathrm{fin}}(\varepsilon)/P\), and the \((\mathrm{RCRT}_\varepsilon)\) threshold approximately closes as the self-consistent bisection
  \[
  \boxed{\;
  P^2\ \gtrsim\ 2\,B_{m,q}^{\mathrm{fin}}(\varepsilon)\;}
  \qquad
  \Longleftrightarrow
  \qquad
  \log P\ \gtrsim\ \frac{\log 2+\log B_{m,q}^{\mathrm{fin}}(\varepsilon)}{2}.
  \]
  Furthermore, if we use the certificate envelope \(\gamma_m^{\mathrm{cert}}(\varepsilon)\) from \S\ref{subsec:pom-quantile-budget-law} (see \(\gamma_m^{\mathrm{cert}}(\varepsilon)=\inf_{q\ge 2}\frac{\log r_q-\log 2+\frac{1}{m}\log(1/\varepsilon)}{q-1}\)) and write \(B_m^{\mathrm{cert}}(\varepsilon):=\exp(m\gamma_m^{\mathrm{cert}}(\varepsilon))\), then this square-root frontier can be written as
  \[
  \boxed{\;
  \log P\ \gtrsim\ \frac{\log 2+m\,\gamma_m^{\mathrm{cert}}(\varepsilon)}{2}
  \qquad(\text{certified square-root quota law}).\;}
  \]
\end{itemize}
\end{corollary}

\subsubsection{Localization of anomalies: intra-fiber Fourier variation and invertible reconstruction}\label{subsec:pom-anom-localization-fourier}
\begin{definition}[Intra-fiber character average (Fourier coefficient)]\label{def:pom-fiber-fourier-coeff}
Fix resolution \(m\) and a finite abelian lift group \(G\), and view the lift label as a function
$$
\ell:\Omega_m\to G.
$$
For any character \(\chi\in\widehat G\), define the intra-fiber character average
$$
\mu_\chi(x):=\frac{1}{d_m(x)}\sum_{\Fold_m(a)=x}\chi(\ell(a)),
\qquad x\in X_m.
$$
\end{definition}

\begin{proposition}[Character average as conditional expectation; residual energy as conditional variance]\label{prop:pom-fiber-fourier-variance}
Let \(f_\chi:=\chi\circ \ell\).
For any probability distribution \(\mu\in\Delta(\Omega_m)\), denote its pushforward \(\pi=(\Fold_m)_*\mu\in\Delta(X_m)\), and denote
\[
E_{\mu,m}[\,\cdot\mid\Fold_m]:L^2(\Omega_m,\mu)\to L^2(X_m,\pi)
\]
the conditional expectation operator (fiber conditional average over \(x\in X_m\)).
Then:
\begin{enumerate}
  \item If \(\mu\) is conditionally uniform on each fiber (e.g., \(\mu=Q_m\pi\) or \(\mu\) is the uniform distribution on \(\Omega_m\)), then for each \(x\in X_m\),
  \[
  (E_{\mu,m}[f_\chi\mid\Fold_m])(x)=\mu_\chi(x),
  \]
  and this coordinate formula is consistent with the counting average of the folding conditional expectation \(E_m\) (Proposition \ref{prop:fold-conditional-expectation}), hence can also be written as \((E_m f_\chi)(x)=\mu_\chi(x)\).
  \item Under any reference distribution \(\mu\), the exact variance decomposition holds:
  $$
  \|f_\chi-E_{\mu,m}[f_\chi\mid\Fold_m]\|_{2,\mu}^2
  =
  \mathbb{E}_\mu\!\big[\mathrm{Var}_\mu(f_\chi\mid \Fold_m)\big],
  $$
  so the energy of the intra-fiber ``non-constant part'' is precisely the average of the conditional variance under the fiber conditional distribution (numerical verification in Table \ref{tab:fold_conditional_variance_decomposition}).
\end{enumerate}
\end{proposition}
\input{sections/generated/tab_fold_conditional_variance_decomposition}

\begin{corollary}[Invertible reconstruction of the intra-fiber label distribution]\label{cor:pom-fiber-label-distribution-inversion}
For each \(x\in X_m\), define the intra-fiber label distribution
$$
p_x(g):=\frac{1}{d_m(x)}\bigl|\{a\in\Omega_m:\ \Fold_m(a)=x,\ \ell(a)=g\}\bigr|.
$$
Then \(\mu_\chi(x)=\sum_{g\in G}p_x(g)\chi(g)\) is the Fourier transform of \(p_x\); hence for finite abelian groups there is the explicit inversion
$$
\boxed{\;
p_x(g)=\frac{1}{|G|}\sum_{\chi\in\widehat G}\mu_\chi(x)\,\overline{\chi(g)}\; }.
$$
Therefore ``anomaly/mixing'' can not only be detected by the global finite-part signature \(\mathsf{Anom}_G\) (Definition \ref{def:pom-anom-signature}), but can also be localized within fibers as invertible local statistics: which labels mix within the same fiber, and the intensity and structure of mixing, are all determined by the local spectrum \(\{\mu_\chi(x)\}_{\chi}\).
\end{corollary}

\begin{proposition}[Fourier distribution law: strict character decomposition of the \(q\)-moment after lift]\label{prop:pom-mom-lift-fourier-distribution}
Following the lift label \(\ell:\Omega_m\to G\) of Definition \ref{def:pom-fiber-fourier-coeff}.
For each \(x\in X_m\) and \(g\in G\), denote
\[
d_m(x,g):=\bigl|\{a\in\Omega_m:\ \Fold_m(a)=x,\ \ell(a)=g\}\bigr|,
\qquad
d_m(x)=\sum_{g\in G}d_m(x,g).
\]
For each character \(\chi\in\widehat G\), define the Fourier transform of the fiber count
\[
\widehat d_m(x,\chi):=\sum_{g\in G} d_m(x,g)\,\overline{\chi(g)}.
\]
Then for any integer \(q\ge 2\) the strict identity holds:
\[
\boxed{\;
\sum_{x\in X_m}\sum_{g\in G} d_m(x,g)^q
\;=\;
\frac{1}{|G|^{q-1}}
\sum_{\substack{\chi_1,\dots,\chi_q\in\widehat G\\ \chi_1\cdots\chi_q=\chi_0}}
\ \sum_{x\in X_m}\ \prod_{i=1}^q \widehat d_m(x,\chi_i)\;}
\]
where \(\chi_0\) is the trivial character.
\end{proposition}
\begin{proof}
For each fixed \(x\), by Fourier inversion
\(
d_m(x,g)=\frac{1}{|G|}\sum_{\chi\in\widehat G}\widehat d_m(x,\chi)\chi(g)
\).
Substituting into \(d_m(x,g)^q\) and summing over \(g\in G\), expanding and using character orthogonality
\(
\sum_{g\in G}\chi_1(g)\cdots\chi_q(g)
=|G|\cdot \mathbf{1}_{\chi_1\cdots\chi_q=\chi_0}
\)
to sieve out the constraint \(\chi_1\cdots\chi_q=\chi_0\) yields the stated identity.
\end{proof}

\begin{remark}[From global signature to local debugging patch]\label{rem:pom-anom-localization-bridge}
The \(\mathsf{Anom}_G(K;\theta)\) of Definition \ref{def:pom-anom-signature} is a constant-level global signature of ``commutativity failure,'' suited for composable verification of rewriting certificates and completion (Remark \ref{rem:pom-anom-completion}).
Corollary \ref{cor:pom-fiber-label-distribution-inversion} in turn provides a finer-grained local syntactic interface: when significant non-singleton behavior of \(p_x\) is observed on certain fibers \(x\), the ``lift channels needing splitting/relabeling'' can be made concrete as a local rewrite of \(\ell\) (splitting the mixed label clusters into independent buckets), thereby upgrading the anomaly from a ``signature'' to a ``structure localization quantity capable of generating patches.''
\end{remark}

\subsubsection{Functorial readout of the joint phase diagram: endpoints and duality of the two-dimensional pressure surface \texorpdfstring{$P(q,u)$}{P(q,u)}}\label{subsec:pom-joint-phase-functor}
The two-dimensional pressure surface \(P(q,u)\) has already been constructed and reproducibly computed in \S\ref{subsec:pom-pressure-surface-qu} via the weighted collision moment-kernel.
Its endpoints and duality operations in two directions carry functorial significance in the calculus:
\begin{itemize}
  \item \textbf{(\(q\)-endpoint)} Fixing \(u\), the map \(q\mapsto P(q,u)\) via the Legendre--Fenchel transform yields the ``space loss multifractal'' upper envelope (the isomorphic operation of Proposition \ref{prop:pom-fold-multifractal-envelope}; here \(r_q\) is simply replaced by \(r_q(u)\));
  \item \textbf{(\(u\)-endpoint)} Fixing \(q\), the tropicalization slope at \(u\downarrow 0\),
  \(
  \lim_{u\downarrow 0}\frac{\log r_q(u)}{\log u}
  \),
  yields the minimum average energy certificate after \(q\)-fold collision compilation (degenerating to Proposition \ref{prop:pom-proj-tropical-minmean} for \(q=1\)).
\end{itemize}
Therefore \(P(q,u)\) can be viewed as the joint semantic object of the same projection word along the ``space moment spectrum--time temperature tilt'' two axes: it locks the certified rewriting and cost-functorial readout of \(\mathbf{PW}\) onto the same computable phase diagram.
